\documentclass[10pt]{article} % For LaTeX2e
\usepackage{tmlr}
% If accepted, instead use the following line for the camera-ready submission:
%\usepackage[accepted]{tmlr}
% To de-anonymize and remove mentions to TMLR (for example for posting to preprint servers), instead use the following:
%\usepackage{tmlr}

\input{math_commands.tex}

\usepackage{hyperref}
% Without an explicit setup hyperref gives every link a 1pt border, which readers see as red
% boxes around each \ref and \cite (124 of them in this build). Hide them.
\hypersetup{hidelinks}
\usepackage{url}
\usepackage{booktabs}
\usepackage{amsmath}
\usepackage{amssymb}
\usepackage{graphicx}

\title{From Controllable Reasoning Effort to Controllable Agent Effort:
Joint Effort Allocation for Tool-Using Agents}
% TMLR G3 anonymous: author block redacted for blind review; restore on acceptance (whitelist: no identity in source)


% Authors are hidden in the anonymous submission (TMLR G3).

\newcommand{\fix}{\marginpar{FIX}}
\newcommand{\new}{\marginpar{NEW}}

\newcommand{\eff}{b}
\newcommand{\reff}{\eff_{\mathrm{reason}}}
\newcommand{\stepeff}{\eff_{\mathrm{steps}}}
\newcommand{\tooleff}{\eff_{\mathrm{tools}}}
\newcommand{\searcheff}{\eff_{\mathrm{search}}}
\newcommand{\viseff}{\eff_{\mathrm{vision}}}

\begin{document}

\maketitle

\begin{abstract}% CANDIDATE-A -- ready, still needs apply-GO. This replaces the 05:19Z draft, whose own
% header forbade applying it because its clause quoted the record-level statistic that C51
% shows to be pseudo-replicated. The clause now carries the unit-level reading instead: 16
% independent units, 3 worse and 0 better, p=0.25, i.e. not significant. Carrier:
% evidence/H4_PSEUDOREPLICATION.txt, row `1_low->4_max`.
% Paste length is re-measured by check_h4p_wording.py against the 1920-character ceiling.
% The superseded draft is kept unmodified beside this file for the audit trail.
% 2026-09-25 (second-scale round): clause (v) now names BOTH registered arms (8B and 14B) and
% says the 14B arm contradicts the 8B mechanism; "that assumption fails" was downgraded to
% "the assumed shape does not hold at either", and the unit-level p=0.25 is labelled "not
% established". Compressions in (ii)-(iv) are wording-only, no number was dropped. Paste
% re-measured at 1907 chars (ceiling 1920) by abstract_readiness_census.py.
Agents spend several resources at once: reasoning tokens, trajectory
length, tool calls, search calls, visual probes. The V4.1-Flash technical report
controls one by reinforcement learning: a scalar effort level
$b\in\{1,\dots,100\}$ whose reward penalises tokens exponentially. We route the other four jointly under one
scalar-budget accounting. (i)~Re-anchoring the
premise: $b$ acts on reasoning tokens only, and only the report's own checkpoints
instantiate its curves; we neither reproduce nor test them. (ii)~On our host we
measure $18$ tool primitives ($\geq30$ samples each; per-primitive
$n$ and a failed probe: Sec.~\ref{sec:e1}): a local call costs
$7.09$~ms against $2233.75$~ms for internet search, a $315\times$ gap; vision is
costliest per unit at $177.34$~ms. (iii)~With those costs, a source-anchored
reasoning kernel and four declared agent-axis gains, the frontier of an \emph{idealised
rational allocator} over $8{,}820$ allocations needs $3.44\times$, $4.52\times$, $2.05\times$ fewer tokens at matched success than
fixed-mid, reasoning-max, agent-max, conditional on the assumed gains; decoupled per-axis budgets cost
$1.32\times$ more at $40\%$ success. The lesson is re-allocation, not more spend; the mix
inverts with budget. (iv)~An end-to-end loop on auto-graded tasks gives monotone
saturating returns: $0.500$ at the low budget, $1.000$ at every larger one, $16$
units per level, no gain from $4\times$ that budget. (v)~With a real language model in the loop ($672$ episodes/arm, two scales) the assumed
shape does not hold at either: $8$B peaks at $0.9375$ on the low budget and falls to $0.7500$
at $8\times$ that budget ($16$ units, McNemar $p{=}0.25$: not established), verification
collapsing at max ($31$ rewrites, no test); $14$B peaks at mid and its tests per episode
\emph{rise} ($1.19\to2.00$), contradicting that mechanism. No gain is identifiable at either
scale (CI widths $1.00$/$1.00$/$0.58$). What the paper contributes is the audited distance
between (iii) and (v): the frontier is a normative benchmark, not a description of any
deployed model, and neither scale's turnover is established on $16$ units. The $8$B curve reproduces across stacks
(3090/vLLM 0.27.1 to 4080/0.30.0): success unchanged in all $23$ cells. Re-derivation: \texttt{analysis.py},
\texttt{h4\_derived\_stats.py}; raw JSON released.
\end{abstract}

\section{Introduction}

An agent solving a long-horizon task spends several resources at once. It spends
\emph{reasoning tokens} inside a turn, \emph{turns} on the trajectory,
\emph{local tool calls} (shell, patch, test runner), \emph{network search calls},
and \emph{visual probes} when the task carries screenshots. In deployment these
axes compete for one budget, and they are paid for in different currencies:
tokens for the model, wall-clock for the tools.

The V4.1-Flash technical report~\citep{v41flash2026} makes one of these axes
controllable. It introduces a scalar effort level as an explicit conditioning
signal during reinforcement learning, prepended to the system prompt:
\begin{quote}
``We therefore introduce a scalar effort level $\eff$ as an explicit
conditioning signal during reinforcement-learning. This mechanism is applied to
both single-turn reasoning and multi-turn agentic tasks.'' (lines 2042--2044)
\end{quote}
\begin{quote}
``Here, $\eff \in \{1, \ldots, 100\}$ denotes the requested effort level.''
(line 2051)
\end{quote}
Training samples $M_b$ responses per prompt per effort level and mean-centers
rewards \emph{within} each $(\text{prompt}, \eff)$ subgroup, so ``responses from
different effort levels are not directly compared'' (line 2063). Effort-dependent
behaviour is induced by a length penalty whose coefficient decays exponentially
with the requested effort (lines 2081--2096):
\begin{equation}
r^{\mathrm{len}}(\ell, \eff) = -\min\!\Big( C_{\max},\; k(\eff)\frac{\ell}{L_{\mathrm{norm}}} \Big),
\qquad
k(\eff) = k_0 \exp\!\Big(-\frac{\eff - \eff_{\min}}{\tau}\Big),
\label{eq:k}
\end{equation}
with $\tau = \lambda\Delta_b$ (line 2089), $k_0$ the penalty coefficient at the
lowest effort level, and $\Delta_b$ the average spacing between training effort
levels (lines 2092--2096). Appendix~C repeats the same coefficient as Eq.~(12)
(lines 3501--3511) and derives a local affine trend between requested effort and
preferred reasoning length under an exponentially decaying marginal benefit
(lines 3743--3773), while explicitly warning that the derivation ``is a local
reward-level approximation, not a claim that measured average lengths must be
linear or pointwise monotonic'' (lines 3779--3780).

The report also shows (Sec.~\ref{sec:related} places this next to the budgeted-agent literature) that this control transfers to agentic use: raising effort
from $25$ to $100$ improves average Pass@1 on eight reasoning-intensive
benchmarks from $67.1\%$ to $76.3\%$, on DeepSWE~v1.1 from $66.0\%$ to $74.2\%$,
and on Terminal-Bench~2.1 from $82.4\%$ to $90.6\%$, at roughly $2.5\times$ more
output tokens (lines 2434--2436), and the gains are front-loaded---the $60$--$80$
range already recovers most of the maximum accuracy at less than half its token
budget, whereas the final step to effort $100$ lengthens agent trajectories by
$1.6$--$1.8\times$ for marginal improvement on Terminal-Bench~2.1
(lines 2439--2441).

\paragraph{The gap.}
\eff{} is a scalar over \emph{reasoning tokens}. The other four axes remain
uncontrolled: the report fixes the scaffold and varies the effort tier inside it
(lines 2456--2465), and its multi-agent analysis scales a wall-clock deadline
rather than an effort budget (lines 2736--2763). Conversely, the agent-economics
literature prices actions and budgets them
\citep{wu2026ecoagentbench,liu2025budgettooluse}, and the test-time-compute
literature allocates compute within reasoning
\citep{zhou2026coba,yang2026gambit,alomrani2025reasoningbudget}, but neither side
places reasoning tokens, trajectory steps, local tools, search calls and visual
probes in one accounting unit with one budget. No prior work we are aware of
routes \emph{all five} axes jointly under a single scalar budget, nor reports how
the optimal \emph{mix} moves as that budget changes.

\paragraph{Contributions.}
\begin{enumerate}
\item \textbf{Re-anchoring and promotion surface} (Sec.~\ref{sec:method}). We
  quote the source definition of \eff{} and $k(\eff)$ with line numbers, and we
  state the promotion surface $\{\reff, \stepeff, \tooleff, \searcheff, \viseff\}$ explicitly. We
  also state the boundary of what the source can support: only the source's own
  checkpoints instantiate the reported curves, and its weights are not available
  to us, so every statement about the source model's controllability in this
  paper is a quotation of the report, never a measurement of ours.
\item \textbf{A measured cost model} (Sec.~\ref{sec:e1}). We profile $18$ tool
  primitives on the host that runs our agent loop ($\geq 30$ samples each,
  per-invocation wall clock, failures reported rather than dropped). The
  measured spread is the reason a joint router has anything to do: local tool
  calls cost $7.09$~ms, a vision unit costs $177.34$~ms, and an internet search
  call costs $2233.75$~ms.
\item \textbf{A joint allocation frontier as a normative benchmark} (Sec.~\ref{sec:e2}). On $8{,}820$
  candidate allocations we compute the cost--success frontier implied by the idealised
  allocator of Sec.~\ref{sec:model}, and compare it
  against max-everything, fixed-mid, single-axis-max and decoupled-budget
  routers, with a $200$-replicate bootstrap over the measured latencies and a
  $30$-case sensitivity grid. One baseline \emph{dominates} the other three.
\item \textbf{An end-to-end check} (Sec.~\ref{sec:e3}). A real agent loop with a
  hard per-episode budget, real tool calls and hidden auto-graded tests shows
  monotone, saturating returns and a budget floor below which the task cannot be
  finished at all.
\item \textbf{The measured gap between that benchmark and real closed-loop behaviour, at two
  scales}
  (Sec.~\ref{sec:res-h4}).
  The loop of the previous item is re-run with a real language model instead of
  the rule-based repairer, at $8$B and, as the protocol registered, at $14$B. The
  saturating-gain assumption that the allocation study rests on is contradicted by
  both arms, which do not contradict it in the same way. On the $8$B arm the
  point-estimate curve peaks at the low budget and then falls, and that fall is
  reached by counting every record separately; on the $16$ independent units the same
  contrast is $3$ worse and $0$ better, $p{=}0.25$, so the decline is observed but not
  established. The $14$B arm does not reproduce it: its curve peaks at the mid budget,
  its low-to-max contrast has more improvements than worsenings, and its test
  invocations per episode rise with the budget rather than falling. What both arms do
  support is the smaller claim: the fitted gains are not identifiable at either scale,
  and neither measured curve has the assumed shape.
\end{enumerate}

\paragraph{Honest scope.}
Our experiments are \emph{not} run on the source model; the protocol for each is in Sec.~\ref{sec:experiments}. All of them are run on an
x86-64 CPU host (96 vCPUs, 754~GB RAM); Sec.~\ref{sec:e1}, Sec.~\ref{sec:e2} and
Sec.~\ref{sec:e3} use a rule-based agent policy, and the joint-allocation study is
a model-based optimisation whose cost side is measured and whose success side is
anchored to published numbers plus four declared assumptions.
Section~\ref{sec:res-h4} additionally runs the loop with a real language model on
a GPU, and it is the one place where the gain assumption is tested rather than
assumed. Section~\ref{sec:limits} lists every assumption and what would
falsify it.

\section{Related Work}
\label{sec:related}

\paragraph{Controllable test-time compute.}
Budget-aware test-time compute is surveyed by \citet{alomrani2025reasoningbudget}.
Within a single reasoning episode, \citet{zhou2026coba} route the next unit of
compute among generation, verification and stopping, reporting $85.13\%$ macro
accuracy while using $49.1\%$ fewer parameter-weighted tokens than a
self-evaluation weighted-voting proxy at $85.20\%$; \citet{yang2026gambit}
concentrate compute on promising partial trajectories by thought-level beam
search, reporting up to $+6.7$ absolute accuracy on HMMT-24 and a $68.5\%$ token
reduction against parallel sampling. \citet{wang2023selfconsistency} sample and
vote. All of these allocate \emph{within} a reasoning episode and treat the
sampling/voting/verifying split as the decision variable. None of them makes
trajectory length, tool calls, search depth or visual probing part of the same
decision, which is the axis set an agent actually pays for; the long-horizon task setting itself follows the SWE-bench line of work \citep{jimenez2024swebench}.

\paragraph{Budgeted agents.}
\citet{wu2026ecoagentbench} build $304$ real-derived tasks in which every task
specifies priced actions and an explicit budget, and test four decisions
(avoid escalation, escalate when local evidence is insufficient, select a model
tier, stop on unsupported premises). They report that tool-API agents reach only
$3.9$--$24.0\%$ micro strict success (at most $7.3\%$ economic consistency), and
that a threshold-crossing budget sweep moves one model's escalation rate from
$0\%$ to only $3\%$---evidence that completion under a budget and economical
action selection are distinct properties. \citet{liu2025budgettooluse} add a
budget tracker so a tool-using agent can observe its remaining budget. Both make
the agent \emph{aware} of cost; neither derives the optimal composition of a
budget across heterogeneous axes, which is our question. Our
Sec.~\ref{sec:e2} result that the optimal mix rotates across the budget is
consistent with their finding that awareness alone does not produce economical
selection.

\paragraph{Source report.}
\citet{v41flash2026} is a systems report about KV-cache compression; its
effort-control section (lines 2039--2109) and appendix~C (lines 3486--3511) are
the only place where a scalar effort interface is defined and calibrated. We use
exactly those two locations, and we use them as \emph{source}, not as a baseline
we run.

\section{Method}
\label{sec:method}

\subsection{The five effort axes and their cost currencies}

The model of Sec.~\ref{sec:model} is instantiated below. We write an allocation as a tuple
$a = (\reff, \stepeff, \tooleff, \searcheff, \viseff)$, where $\reff$ is the requested reasoning
effort on the source's own $1$--$100$ scale, $\stepeff$ is the trajectory length in
turns, and $\tooleff$, $\searcheff$, $\viseff$ count local tool calls, network search calls and
visual probes. Two costs are tracked for every allocation. The \emph{token} cost
\begin{equation}
\mathrm{tok}(a) = \stepeff L(\reff)
  + \stepeff c_{\mathrm{step}}
  + (\tooleff + \searcheff) c_{\mathrm{tool}}
  + \searcheff c_{\mathrm{extra}}
  + \viseff c_{\mathrm{vis}},
\end{equation}
where $L(\reff)$ is the reasoning tokens per turn, is the primary currency
because it is the quantity the source's RL penalises \eqref{eq:k}. The
\emph{latency} cost is
\begin{equation}
\mathrm{ms}(a) = \stepeff \frac{L(\reff)}{1000} r_{\mathrm{dec}}
  + \stepeff \tau_{\mathrm{step}}
  + \tooleff \tau_{\mathrm{tool}}
  + \searcheff \tau_{\mathrm{search}} \big(1 + (1-\pi_{\mathrm{search}})(\rho-1)\big)
  + \viseff \tau_{\mathrm{vis}},
\end{equation}
where every $\tau$ is a measured per-invocation latency (Sec.~\ref{sec:e1}) and
$\pi_{\mathrm{search}}$ is the probability that a search call returns a usable
result, with wasted cost $\rho$ on a failed call. Pricing tokens at fixed rates
plus measured tool time follows the report's own latency accounting, which
``assign[s] costs from token counts at fixed prefill/decode rates plus measured
tool-execution time'' (lines 2742--2743). The token prices
$c_{\mathrm{step}}=1000$, $c_{\mathrm{tool}}=200$ and $c_{\mathrm{vis}}=200$
and the decode rate $r_{\mathrm{dec}}=30$ per $10^3$ tokens are
\emph{assumptions}, declared in Sec.~\ref{sec:limits} and swept in
Sec.~\ref{sec:e2}.

\subsection{Source-anchored reasoning kernel}
\label{sec:kernel}

We parameterise the success contribution of the reasoning axis directly from the
report's numbers (this kernel is calibrated in App.~\ref{app:calibration}). Let $w$ be a concave map of effort onto the $25 \to 100$
accuracy swing with $w(25)=0$ and $w(100)=1$,
\begin{equation}
w(b) = \frac{\log\!\big(1+\tfrac{b-b_{\min}}{s_0}\big) - \log\!\big(1+\tfrac{25-b_{\min}}{s_0}\big)}
             {\log\!\big(1+\tfrac{100-b_{\min}}{s_0}\big) - \log\!\big(1+\tfrac{25-b_{\min}}{s_0}\big)},
\end{equation}
so that
$\mathrm{acc}(\reff) = 0.671 + (0.763-0.671) w(\reff)$ reproduces the published
endpoints $67.1\%$ and $76.3\%$ (line 2434) for any $s_0$. We set $s_0=5000$ by
solving the report's front-load statement (lines 2439--2440): the fraction of the
$25 \to 100$ accuracy gain recovered at $\reff=80$ is $0.735$ for this $s_0$.
For the length model we use $L(\reff) = L(25)\exp(\kappa(\reff-25))$ with
$\kappa = \log(2.5)/(100-25) = 0.01222$, which reproduces the report's
``roughly $2.5\times$ more output tokens'' (line 2436) exactly at the endpoints
and matches the affine effort--length trend of appendix~C Eq.~(16)
(lines 3756--3767) in the small-$\kappa$ regime.

\subsection{Joint model}
\label{sec:model}

Given the kernel of Sec.~\ref{sec:kernel} and the hypotheses of Sec.~\ref{sec:pre-registration}, the agent axes are modelled as \emph{substitutable} routes that each recover part
of the residual failure probability, with exponential saturation:
\begin{equation}
\mathrm{agent}(a) = 1 - \prod_{x \in \{\stepeff,\tooleff,\searcheff,\viseff\}}
   \Big(1 - g_x\big(1 - e^{-x/\mathrm{ref}_x}\big)\Big),
\end{equation}
\begin{equation}
P(\text{solve} \mid a) = 1 - \big(1 - \mathrm{acc}(\reff)\mathrm{agent}(a)\big)(1-p_0),
\label{eq:joint}
\end{equation}
with $p_0 = 0.35$ the base success probability. Eq.~\eqref{eq:joint} is monotone
non-decreasing in every axis and bounded by $1$; the four gains $g_x$ and their
half-saturation references $\mathrm{ref}_x$ are \emph{assumptions}
(Sec.~\ref{sec:limits}). The substitutability is the modelling choice that makes
the problem non-trivial: with a conjunctive model (every axis must succeed)
almost every allocation fails and the frontier collapses.

\textbf{What this model is for.} We call the saturating form above together with
Eq.~\eqref{eq:joint} the \emph{idealised rational allocator}. It is the target a
budget-limited agent \emph{ought} to solve, and it is not offered as a description of any
deployed language model: the gains $g_x$ are declared assumptions (Sec.~\ref{sec:limits}),
and by construction the form is monotone in every axis, so it cannot represent an allocation
that does worse with more budget. The empirical object of this paper is the \emph{distance}
between that benchmark and what two scales of a real model do under the same budgets
(Sec.~\ref{sec:res-h4}); the frontier is the measuring stick, not the claim.

\subsection{Pre-registered hypotheses}
\label{sec:pre-registration}

We fixed the following before running the allocation study. \textbf{H1}: at
matched cost, the joint router attains higher success than a fixed mid setting on
every axis. \emph{Falsified if} the success difference is $\leq 0$ at the
baseline's own cost. \textbf{H2}: allowing budget to be traded freely across the
five axes is cheaper than hard-partitioning the budget into five private shares.
\emph{Falsified if} the decoupled partition reaches a given success level at equal
or lower cost. \textbf{H3}: the composition of the optimal allocation changes
across the budget rather than scaling uniformly. \emph{Falsified if} the optimal
levels are (nearly) co-monotone in the budget.

\subsection{Declared limits}
\label{sec:limits}

\begin{enumerate}
\item \textbf{No source-model access.} The report states that model checkpoints
  ``are available at'' a model hub (line 24). We did not obtain or run them: the
  weights are not part of this study, and no result here is a measurement of the
  source model. Every controllability statement about it is a quotation with a
  line number.
\item \textbf{Assumed parameters} (declared, swept in the $30$-case grid of
  Sec.~\ref{sec:e2}): the four agent-axis gains
  $g_{\mathrm{steps}}{=}0.30$, $g_{\mathrm{tools}}{=}0.25$,
  $g_{\mathrm{search}}{=}0.28$, $g_{\mathrm{vision}}{=}0.20$; half-saturations
  $\mathrm{ref}{=}(10,6,4,4)$; base success $p_0{=}0.35$; token prices
  $c_{\mathrm{step}}{=}10^3$, $c_{\mathrm{tool}}{=}200$,
  $c_{\mathrm{vision}}{=}200$; decode rate $r_{\mathrm{dec}}{=}30$ per $10^3$
  tokens; search usability $\pi_{\mathrm{search}}{=}0.85$ with retry cost
  $\rho{=}1.4$. These are \emph{not} measured and \emph{not} published by the
  source. Sec.~\ref{sec:res-h4} is an attempt to measure three of them on a real
  model at two registered scales; it did not succeed in pinning them at either
  (the bootstrap intervals are too wide), and its point estimates contradict the
  assumed monotone shape on that task set at both scales, in different directions
  ($8$B falls after the low budget, $14$B after the mid one), so these values
  remain assumptions and the frontier remains conditional.
\item \textbf{Synthetic policy in Sec.~\ref{sec:e3}.} The end-to-end loop uses a
  deterministic rule-based repairer, not a language model. It therefore supports
  claims about loop-level cost accounting and budget exhaustion, not claims about
  any model's insight-per-effort curve. Sec.~\ref{sec:res-h4} supplies the
  model-backed counterpart: its result is a \emph{negative} one at the $8$B scale and
  a \emph{contradictory} one at $14$B, so the pair supports neither a general decline
  nor a general flatness.
\item \textbf{One host.} All latencies come from one container (96 vCPUs,
  754~GB RAM). Absolute latencies are host-specific; the ratios are the portable
  quantity, and Sec.~\ref{sec:e1} reports both.
\item \textbf{Two models, one task set} (Sec.~\ref{sec:res-h4}). The real-model
  evidence is \emph{two} instruction-tuned models of one family---Qwen3-8B-AWQ and
  Qwen3-14B-AWQ, $672$ episodes each, run under identical code and seed on the same
  eight tabular-repair tasks used throughout---so it is a measurement of those two
  pairings and not a general property of language-model agents. The two pairings do
  not agree with each other: the $8$B curve peaks at the low budget and falls, the
  $14$B curve peaks at the mid budget and its paired contrast has more improvements
  than worsenings, and only the $8$B arm defers verification as the budget grows. Two
  sizes are two points, not a scale trend: nothing here separates model size from the
  pairing artefacts, and no third size was run. The tasks are short and
  auto-graded, the hidden suites are fixed, and success is a step function of
  whether the schedule is completed; a longer-horizon task set would move both the
  peak and the point at which behaviour degrades.
\item \textbf{The assumed gain model is monotone.} $P(\text{solve}\mid a)$ in
  Eq.~\eqref{eq:joint} increases in every axis and cannot represent a decrease with
  budget. On the real-model runs the measured curve decreases after the low budget at
  $8$B and after the mid budget at $14$B, so the form does not hold at either scale;
  this is why the fit reaches only $r^2{=}0.291$ on the $8$B cells
  and why we withhold all three fitted gains at both scales. The frontier of
  Sec.~\ref{sec:e2} is
  therefore a \emph{conditional} model, not a description of the measured agent.
\item \textbf{Phase-one artefacts were excluded, and the exclusion is recorded.}
  The real-model run first lost $10$ of $40$ pilot episodes to action-format
  truncation caused by the model's default deliberation mode; that mode was disabled
  before the reported run, whose parsing failures are zero. Two further hazards are
  recorded with their evidence: a co-tenant job used $\approx\!13$~GB of the device
  without taking the shared lock, so the reported wall-clock is an upper bound under
  contention, and the allocator was given free rather than total device memory.
\item \textbf{Accelerator use.} Sec.~\ref{sec:e1}, Sec.~\ref{sec:e2} and
  Sec.~\ref{sec:e3} are CPU-side. Sec.~\ref{sec:res-h4} ran on two accelerators under
  the shared lock: the first $8$B arm on an \texttt{RTX 3090, 49152~MiB}, and---on the
  same harness code---the second $8$B arm and the $14$B arm on an \texttt{RTX 4080,
  32760~MiB}. The device state at load and the run timing are recorded in each arm's
  result file. An earlier draft of this list said no accelerator was used;
  that applied to Rounds~1--2 only and is not true of the reported run.
\end{enumerate}

\section{Experiments}
\label{sec:experiments}

\subsection{E1/E1b: measured tool cost primitives}
\label{sec:e1}

\textbf{What is measured.} We time $18$ primitives that a coding or research
agent actually invokes, each with its own repetition count (all $\geq 30$):
process spawn, shell command, interpreter start-up, small and medium file reads,
atomic file write, directory scan, a real internet round trip, an unreachable
host, a CPU kernel, a four-worker process fan-out, a read-tool call with
correctness verification, PNG encode/decode of a $1280\times720$ frame, a screen
grab attempt, a small test-runner invocation, a loopback HTTP search over $500$
documents, a patch application, and a JSON-emitting CLI call. A tool call counts
as a \emph{semantic failure} when it returns the wrong answer, not merely when it
exits non-zero; the read-tool task deliberately mixes a matching query (expected
$200$ hits) with a non-matching query (expected $0$) so that a confidently wrong
tool is scored as a failure.

\textbf{Reproducibility.} Generator scripts
\texttt{e1\_tool\_latency.py} and \texttt{e1b\_agent\_primitives.py} (under
\texttt{experiments/}); raw per-invocation samples in
\texttt{e1\_tool\_latency.json} and \texttt{e1b\_agent\_primitives.json} (under
\texttt{results/}); aggregate means recomputed by \texttt{analysis.py}, which
asserts agreement with the stored aggregates.

\subsection{E2: joint allocation}
\label{sec:e2}

\textbf{What is computed.} We enumerate $8{,}820$ candidate allocations over
$\reff\in\{0,25,50,75,100\}$, $\stepeff\in\{1,2,4,8,12,20,32\}$,
$\tooleff\in\{0,1,2,4,8,16,24\}$, $\searcheff\in\{0,1,2,4,8,16\}$,
$\viseff\in\{0,1,2,4,8,16\}$, evaluate Eq.~\eqref{eq:joint} and both cost functions
for each, and extract the Pareto frontier. Baselines: max-everything
($a_{\max}$), a fixed mid setting $(50,8,4,2,2)$, reasoning-only max, agent-axes
max, and a \emph{decoupled} router that splits the total budget into five private
shares and lets each axis climb greedily inside its own share. We report two
matched comparisons per baseline: maximum success at the baseline's own cost, and
minimum cost at the baseline's own success. Uncertainty comes from a
$200$-replicate bootstrap that resamples per-invocation latencies from the E1/E1b
raw samples; a $30$-case grid (three values for each of ten parameters) checks
that the ordering is not an artefact of the assumptions.

\subsection{E3: end-to-end loop under budget degradation}
\label{sec:e3}

\textbf{What is executed.} Eight auto-graded program-repair tasks, each with a
deliberately incomplete starter module and a hidden unit-test suite the agent
never sees. A real ReAct-style loop performs actions drawn from a schedule
(\texttt{read}, \texttt{vision-probe}, \texttt{write}, \texttt{test}), where a
\texttt{test} action spawns a real subprocess; exceeding any of the four budgets
(steps, tool calls, context characters, vision probes) terminates the episode as
a failure; success is the subprocess exit code of the hidden suite and nothing
else. Five budget levels $\times$ two schedules $\times$ three repeats
$\times$ eight tasks $= 240$ episodes.

\section{Results}
\label{sec:results}

\subsection{Measured cost primitives}
\label{sec:res-e1}

Sec.~\ref{sec:res-e1} reports the measured cost primitives (Table~\ref{tab:latency}). Three facts drive
everything downstream. First, spawn-dominated actions are not cheap in relative
terms: interpreter start-up is $30.36$~ms and a four-worker fan-out is
$38.34$~ms, an order of magnitude above a file read ($0.033$~ms). Second, the
heterogeneity across tool families is the point: the local read tool costs
$7.09$~ms ($n{=}60$, $95\%$ bootstrap CI $[6.75, 7.41]$) while an internet round
trip costs $2233.75$~ms ($n{=}40$, CI $[1609, 2933]$), a ratio of $315\times$;
a failed network call still costs $961.82$~ms before failing ($n{=}40$, failure
rate $1.000$). Third, the vision axis is the most expensive per unit
($177.34$~ms, CI $[175.02, 181.37]$) and its \emph{capture} step is unavailable
on this headless host: a real screen grab failed in $30$ of $30$ attempts. We report
that as a failure rather than substituting a constant, and we proxy the vision
unit by a frame codec (encode plus decode) instead. The proxy bounds that unit
\emph{from below}: it prices the codec only, not the visual-encoder pass a deployed
multimodal agent pays on the same frame, which we do not measure here. Every vision figure in
this paper is therefore a floor, and the router's preference for buying vision before network
search is conservative in exactly that direction.

The local search alternative is $1.02$~ms ($n{=}60$) against the internet's
$2233.75$~ms, i.e. the same logical action costs $2194\times$ more when it leaves
the host. This is the largest dispersion we measured, and it is why the search
axis stays at zero in the low-latency part of the frontier
(Sec.~\ref{sec:res-e2}).

\begin{table}[t]
\caption{Measured per-invocation cost primitives on the agent host. $n$ is the
number of repetitions, fail is the fraction of invocations that returned a wrong
answer or raised, mean/p50/p95 are over successful invocations for the latency
columns, and the CI is a $4000$-replicate bootstrap of the mean. Units are
milliseconds. Every value is recomputed from \texttt{results/*.json} by
\texttt{analysis.py}.}
\label{tab:latency}
\begin{center}
\small
\begin{tabular}{lrrrrrl}
\toprule
Primitive & $n$ & fail & mean & p50 & p95 & bootstrap CI$_{95}$ of mean \\
\midrule
shell spawn (\texttt{/bin/true})        & 120 & 0.000 & 1.31   & 1.32   & 1.42   & $[1.28, 1.33]$ \\
shell command (\texttt{grep}$+$\texttt{wc}) & 90 & 0.000 & 7.13 & 7.37 & 8.11 & $[6.94, 7.31]$ \\
interpreter start-up                    & 60 & 0.000 & 30.36  & 31.41  & 34.38  & $[29.58, 31.11]$ \\
file read, 4~KiB                        & 200 & 0.00 & 0.033  & 0.034  & 0.037  & $[0.032, 0.035]$ \\
file read, 256~KiB                      & 60 & 0.00 & 0.052  & 0.047  & 0.053  & $[0.048, 0.059]$ \\
atomic file write, 4~KiB                & 90 & 0.000 & 0.51   & 0.45   & 0.96   & $[0.46, 0.57]$ \\
directory scan, 200 files               & 120 & 0.000 & 0.101 & 0.093  & 0.156  & $[0.092, 0.115]$ \\
CPU kernel ($512^2$ matmul)             & 60 & 0.000 & 1.94   & 0.68   & 1.47   & $[0.80, 4.09]$ \\
4-worker process fan-out                & 45 & 0.000 & 38.34  & 37.52  & 47.43  & $[36.32, 40.95]$ \\
read tool, exit \emph{and} answer checked & 60 & 0.000 & 7.09  & 7.61   & 8.70   & $[6.75, 7.41]$ \\
internet round trip (GET, 2~KiB)        & 40 & 0.000 & 2233.75 & 1220.81 & 6693.34 & $[1609, 2933]$ \\
unreachable host (DNS/TLS failure)      & 40 & 1.00 & 961.82 & 101.43 & 5132.05 & $[459, 1583]$ \\
PNG codec, $1280\times720$              & 60 & 0.000 & 177.34 & 174.54 & 186.26 & $[175.02, 181.37]$ \\
screen grab (headless host)             & 30 & 1.00 & ---    & ---    & ---    & unavailable \\
test runner, 3 tests                    & 30 & 0.000 & 58.25  & 58.16  & 63.50  & $[57.05, 59.44]$ \\
loopback HTTP search, 500 docs          & 60 & 0.000 & 1.02   & 0.89   & 1.45   & $[0.82, 1.34]$ \\
patch application                       & 60 & 0.000 & 2.26   & 2.24   & 2.58   & $[2.22, 2.32]$ \\
JSON CLI tool call                      & 60 & 0.000 & 33.98  & 31.75  & 35.74  & $[31.09, 38.22]$ \\
\bottomrule
\end{tabular}
\end{center}
\end{table}

\subsection{Joint allocation: the frontier and the baselines}
\label{sec:res-e2}

The allocation study of Sec.~\ref{sec:e2} gives the headline comparison in Table~\ref{tab:e2}. H1 holds for the three
non-saturating baselines and fails in the literal form we wrote for the
max-everything baseline, which is already at the attainable maximum; we report
both rather than choosing the flattering one.

At the fixed-mid baseline's own token cost ($23{,}429$), the joint router reaches
success $0.6341$ against the baseline's $0.5320$, a gain of $+10.21$ percentage
points. Equivalently, to reach the baseline's success level the router needs
$6{,}810$ tokens, a factor of $3.44$ fewer. Against reasoning-only max the gain
is $+10.16$ points at matched cost and $4.52\times$ fewer tokens at matched
success. Against agent-axes max it is $+2.11$ points and $2.05\times$. Against
max-everything the router reaches the same success ($0.6906$) at the same cost:
once every axis is at its maximum, no allocation can do better in this model, and
spending more is pure waste. What the max-everything baseline \emph{loses} is the
cost of the last few points: to reach $90\%$ of the attainable success
($0.6216$), the router needs $16{,}800$ tokens while max-everything carries
$145{,}600$, a factor of $8.67$ (Table~\ref{tab:e2}).

H2 also holds. With five private budget shares and no cross-axis trade, the
decoupled router needs $3{,}284$ tokens to reach $40\%$ success against the
unified router's $2{,}484$, i.e. $1.32\times$ more expensive, and it tops out at
$0.6736$ against the unified $0.6906$.

\begin{table}[t]
\caption{Joint allocation against four router baselines (token currency).
``gain at baseline cost'' is the success difference at the baseline's own token
cost; ``ratio at matched success'' is baseline cost divided by the router's
minimum cost to reach the baseline's success. The $+0.000$ entry for
max-everything is a real tie, not a rounding artefact: the baseline is at the
attainable maximum.}
\label{tab:e2}
\begin{center}
\small
\begin{tabular}{lrrrrr}
\toprule
Router & success & tokens & gain at & router cost at & ratio at \\
       &         &        & baseline cost & matched success & matched success \\
\midrule
max every axis      & 0.6906 & 145{,}600 & $+0.000$ pts & 145{,}600 & $1.000\times$ \\
fixed mid $(50,8,4,2,2)$ & 0.5320 & 23{,}429 & $+10.21$ pts & 6{,}810 & $3.44\times$ \\
reasoning only max  & 0.5478 & 34{,}400 & $+10.16$ pts & 7{,}610 & $4.52\times$ \\
agent axes max      & 0.6633 & 101{,}717 & $+2.11$ pts & 49{,}600 & $2.05\times$ \\
decoupled budgets   & 0.6736 & 65{,}600 & --- & 3{,}284 (at 0.40) & $1.32\times$ (at 0.40) \\
\bottomrule
\end{tabular}
\end{center}
\end{table}

\paragraph{The frontier is a composition claim, not a spending claim.}
H3 holds, and it is the most useful result. The optimal allocation rotates
completely across the budget (Table~\ref{tab:frontier}). At $5{,}000$ tokens the
optimum is $(25, 1, 4, 2, 4)$: cheap search and moderate vision, almost no
trajectory. At $10{,}000$ tokens reasoning is already at maximum ($100$) while
the trajectory is still a single turn, because extending the trajectory
multiplies the per-turn reasoning cost $L(\reff)$. Only at $20{,}000$ tokens and
above does the router buy turns ($4$, then $8$, then $12$, then $20$), and it
buys them against reasoning that is already saturated. In latency currency the
ordering flips: at a $2{,}000$~ms budget the router refuses any network search
($0$ calls, because one call costs $2233.75$~ms) and buys local tools instead.
An ablation confirms that no single axis substitutes for this rotation. Holding the
other four axes at their mid setting, every single-axis frontier passes through the
same mid point at $0.5320$ (cost $23{,}429$); pushing one axis beyond it buys at most
$0.5827$, and that is \emph{search} alone at $1.4\times$ the mid cost. Reasoning alone
is the weakest of the five frontiers, topping out at $0.5478$ after $1.5\times$ the mid
cost, and reaching the search figure would take steps alone at $3.7\times$. At the mid
cost the joint router reaches $0.634$.

\begin{table}[t]
\caption{Composition of the optimal allocation as the token budget changes. The
mix rotates rather than scaling uniformly, which is H3.}
\label{tab:frontier}
\begin{center}
\small
\begin{tabular}{rrrrrr}
\toprule
token budget & $\reff$ & $\stepeff$ & $\tooleff$ & $\searcheff$ & $\viseff$ \\
\midrule
5{,}000  & 25  & 1  & 4  & 2  & 4 \\
10{,}000 & 100 & 1  & 8  & 4  & 8 \\
20{,}000 & 100 & 1  & 16 & 16 & 16 \\
30{,}000 & 100 & 4  & 24 & 8  & 16 \\
50{,}000 & 100 & 8  & 24 & 16 & 16 \\
80{,}000 & 100 & 12 & 24 & 16 & 16 \\
120{,}000 & 100 & 20 & 24 & 16 & 16 \\
\bottomrule
\end{tabular}
\end{center}
\end{table}

\paragraph{Robustness.}
The $200$-replicate bootstrap over measured latencies gives a success gain of
$0.000$ points with a $95\%$ interval of $[0.000, 0.000]$ for the
max-everything baseline---i.e. the tie is stable, because the baseline saturates
under latency resampling too. The $30$-case sensitivity grid is less
comfortable and we report it in full: the gain over max-everything ranges from
$-7.14$ to $+7.14$ points, with $15$ exact ties (the default regime) and $8$
cases where the router is \emph{worse}. Every losing case is the same mechanism:
when the agent axes are made weak (base success $p_0{=}0.20$, or gains
$0.10$--$0.15$) or their saturation is pushed far out
($\mathrm{ref}_{\mathrm{steps}}{=}16$, $\mathrm{ref}_{\mathrm{tools}}{=}12$),
success becomes accuracy-limited, so the only axis that still moves the frontier
is reasoning, and max-everything's reasoning setting is already at the ceiling.
The three gains over the fixed-mid, reasoning-only and agent-only baselines are
robust across the grid; the tie with max-everything is the fragile claim, and we
do not rely on it.

\subsection{End-to-end loop}
\label{sec:res-e3}

The loop of Sec.~\ref{sec:e3} produced $240$ executed episodes (Table~\ref{tab:e3}). The success curve is
monotone non-decreasing in the budget and saturates early: $0.000$ at the
minimal budget ($2$ steps), $0.500$ at the low budget, and $1.000$ from the mid
budget upward, with no further gain from four times the budget. Cost per success
is $51.6$~ms at the low budget against $52.0$ and $51.0$~ms at the mid and max
budgets: beyond saturation, extra effort is pure cost. This reproduces at loop
level the same shape the source report describes for the reasoning axis
(``the gains are front-loaded'', line 2439).

The information in this curve is concentrated in its lower two levels: three of
the five budget levels (mid, high, max) sit at $1.000$, so the upper half of this
particular frontier carries no discriminating information. The task set used here
saturates at $8$ steps; ``buying more steps yields nothing'' is therefore a
statement about the ceiling of \emph{this eight-task set}, not a general law about
agent effort. A harder task set would move the saturation point and reopen the
upper region, and we flag this as the main way the E3 result under-reports the
regime a real agent operates in.

The schedule decomposition is also informative: at the low budget the
read-then-write schedule succeeds $24/24$ while the vision-probe-heavy schedule
succeeds $0/24$, purely because its first action needs a vision probe that the
budget does not grant. The budgets, not the policy, decide. This is a second
ceiling effect: the schedule is fixed, so measured success is a step function of
whether the whole schedule fits, not a smooth function of budget.

\begin{table}[t]
\caption{End-to-end agent loop under degraded budgets. $48$ records per level = $16$
independent (task, schedule) units $\times$ $3$ deterministic repeats; the policy is
deterministic, so the copies carry no information. Success is the hidden test suite's exit
code; the CI is an exact binomial (Clopper--Pearson) on those $16$ units, printed to 3 dp rounded outward, never a bootstrap over the $48$ rows, which would print a zero-width interval at $16/16$. Cost per success is undefined at the level where nothing succeeds.}
\label{tab:e3}
\begin{center}
\small
\begin{tabular}{lrrrrrl}
\toprule
Budget level & max steps & records ($n$ units) & success & CI$_{95}$ & mean ms & cost/success \\
\midrule
minimal & 2  & 48 (16) & 0.000 & $[0.000, 0.206]$ & 0.014 & --- \\
low     & 4  & 48 (16) & 0.500 & $[0.246, 0.754]$ & 25.8 & 51.6~ms \\
mid     & 8  & 48 (16) & 1.000 & $[0.794, 1.000]$ & 52.0 & 52.0~ms \\
high    & 16 & 48 (16) & 1.000 & $[0.794, 1.000]$ & 51.7 & 51.7~ms \\
max     & 32 & 48 (16) & 1.000 & $[0.794, 1.000]$ & 51.0 & 51.0~ms \\
\bottomrule
\end{tabular}
\end{center}
\end{table}

\subsection{A real-model test of the gain assumption}
\label{sec:res-h4}

The allocation study of Sec.~\ref{sec:e2} \emph{assumes} that each agent axis has a
saturating gain $g_x$. Sections~\ref{sec:res-e1}--\ref{sec:res-e3} hold the policy
fixed, so none of them tests that assumption. We test it by replacing the rule-based
repairer of Sec.~\ref{sec:e3} with a real instruction-tuned model: Qwen3-8B-AWQ served
by vLLM, temperature $0$, one JSON action per turn, over the same eight tasks, the same
five budget levels and the same two action menus. The protocol registers a second, larger
arm (Qwen3-14B-AWQ); it was run, and because it does not agree with the $8$B arm it is
reported in its own paragraph at the end of this subsection rather than in a footnote.
The run has two parts: $240$
episodes replicate Sec.~\ref{sec:e3} exactly ($8$ tasks $\times$ five levels
$\times$ two menus $\times$ three repeats), and $432$ further episodes sweep one
axis at a time while holding the others at the mid allocation, which is what the
gain fit below needs in order to separate the axes. Grading is unchanged---the exit code of the same hidden suites, which the model never sees.

\textbf{What the model is shown (budget awareness).} The budget is explicitly visible in the
model's context, not only enforced outside it. Every turn is built from one template
(\texttt{experiments/h4\_real\_llm\_gx.py}, \texttt{SYS} at line 406 and \texttt{USER\_TMPL}
at lines 419--429) which prints \texttt{budget:} followed by \texttt{max\_steps},
\texttt{max\_tool\_calls}, \texttt{max\_context\_chars} and \texttt{max\_vision\_probes}, then
a \texttt{state} object carrying \texttt{steps\_used}, \texttt{tools\_used},
\texttt{vision\_used} and the action list, then the serialised history, then a request for
the next single JSON action. The model therefore sees both its ceiling and its own
consumption at every step. This fixes the causal reading of the $8$B collapse reported below:
the disappearance of \texttt{test} actions as the budget grows is a decision taken under a
\emph{known} remaining budget---deferral of the terminating action---and the competing
explanation, context drift under a truncation the model cannot see, does not apply to this
harness.

Two defects had to be fixed first, and both would have produced a false measurement.
First, a pilot run lost $10$ of $40$ episodes to action-parsing failures: in its default
reasoning mode the model emits $\approx\!974$ tokens of deliberation before the action
and is truncated before the JSON is reached, which censors the budget--success
relationship for reasons unrelated to budget. Disabling that mode reduced a turn from
$7.46$~s to $0.58$~s and removed the parse failures entirely. Second, the target host was
shared by a co-tenant that did not take the shared GPU lock, so the allocator was given
the \emph{free} device memory rather than a fixed fraction of total; the run records its
wall-clock ($1818.2$~s, summed from per-episode timers) and its device state at load, and
we treat the reported timing as an upper bound under contention.

\begin{table}[t]
\caption{The $8$B arm: a real language model in the loop of Sec.~\ref{sec:e3}. $672$ episodes
over the $8$ tasks, each repeated $3$ times: $240$ replicate the five budget levels
of Sec.~\ref{sec:e3} under both action menus ($8\times5\times2\times3$), and $432$
come from the three single-axis sweeps used for the gain fit ($120$ budget, $168$
tool-call and $144$ vision-probe settings, with the remaining axes held at the mid
allocation). The level column pools the $240$ replication episodes only and reports
their $n$. Success is the hidden suites' exit code. Every failure at every level is
a step-budget exhaustion; action-parsing failures are zero throughout. The last two
columns are the mean number of file rewrites and of test invocations per episode.}
\label{tab:h4}
\begin{center}
\small
\begin{tabular}{lrrrrr}
\toprule
Budget level & max steps & $n$ & success & writes/ep & tests/ep \\
\midrule
minimal & 2  & 48 & 0.0000 & 1.00 & 0.00 \\
low     & 4  & 48 & \textbf{0.9375} & 1.69 & 0.94 \\
mid     & 8  & 48 & 0.8750 & 4.81 & 0.94 \\
high    & 16 & 48 & 0.8125 & 7.38 & 0.88 \\
max     & 32 & 48 & \textbf{0.7500} & 13.81 & 0.75 \\
\bottomrule
\end{tabular}
\end{center}
\end{table}

Table~\ref{tab:h4} reports the $8$B arm, and it does not have the saturating shape. Success
is $0.0000$ at the budget that cannot fit the minimum schedule, peaks at $0.9375$ on the
low budget, and then \emph{declines} monotonically to $0.7500$ at eight times that
budget. Pairing the same tasks and menus across the two extremes gives $9$ record pairs that
worsen and $0$ that improve, an exact McNemar $p{=}0.0039$; that figure counts each record
separately and so is pseudo-replicated, because the three repeats per unit
record the same outcome, which leaves $16$ independent units: on units it is $3$ and $0$, $p{=}0.25$, which does not establish the decline. We therefore state this curve as an observation about these episodes, not as an established property of the model. At $n{=}16$ independent (task, schedule) units, the paired two-sided $\alpha{=}0.05$, power $0.8$ minimum detectable standardized difference is $d_z{=}(t_{.975,15}{+}t_{.8,15})/\sqrt{16}{=}(2.131{+}0.866)/4{\approx}0.75$, so only large effects are detectable and the $3$-vs-$0$ contrast with CI width $1.00$ reflects underpowering, not evidence of absence. All failures at all levels
are step-budget exhaustion, and parsing failures are zero, so the shape is not censoring.

Two claims must be separated here: a model monotone by construction cannot represent a decline, and the unit-level reading establishes no decline. The model of
Eq.~\eqref{eq:joint} is monotone in every axis by construction: $1-p$ is a product of
factors that shrink as any $b_x$ grows, so success can only rise with budget. It cannot
represent the measured decline, which is also why the fit on the $17$ structurally
feasible cells reaches only $r^2{=}0.291$. We therefore report the fitted gains but
withhold them: for all three axes the bootstrap interval is too wide to pin the
parameter (widths $1.00$, $1.00$ and $0.58$ against a $0.30$ threshold), so \emph{no}
$g_x$ value from this run may be quoted, and the ``compute-limited versus
accuracy-limited'' classification that the gains were meant to support is not available
from this experiment.

What the $8$B measurement does show is a mechanism, and it is a behavioural one. The last
two columns of Table~\ref{tab:h4} show that as the advertised budget grows the model
does more of the work that cannot terminate and less of the work that can: rewrites per
episode rise from $1.69$ to $13.81$ while test invocations fall from $0.94$ to $0.75$. At
the largest budget the representative failing trajectory is one read followed by $31$
rewrites and \emph{no test at all}, ending in budget exhaustion. A larger budget is thus
not a larger search over the same behaviour; it changes the behaviour, and it changes it
in the direction of deferring the only action that produces the grading signal. That last
sentence is a claim about the $8$B arm only: the $14$B arm of this subsection's closing
paragraph shows the opposite behaviour on the same task set.

\paragraph{A correction we had to make to ourselves.}
An earlier version of this analysis reported a significant quadratic trend in budget and
read it as an interior optimum. That reading is an artefact. The trend is driven by the
lowest level, where the budget cannot fit ``read, write, test'' and success is $0.0000$
for structural reasons; including it makes the curve look inverted-U. Excluding that
level, the quadratic term is not significant (likelihood-ratio $\chi^2{=}0.19$ on $1$~df,
$p{=}0.66$) and its sign reverses. The defensible statement for the $8$B arm is therefore a
monotone decline above the low budget, and we do not claim an interior optimum \emph{there};
the closing paragraphs of this subsection show that this correction does not carry over to
the $14$B arm, where the same test points the other way. We report this
because the same trap applies to any budget sweep that includes a structurally
infeasible endpoint.

\paragraph{Is the finding an artefact of the machine that first showed it?}
We repeated the whole $672$-episode sweep once more, holding the harness revision, the
model weights, the seed (\texttt{20260911}), the repeat count, the arm set, the chat mode
and the VRAM floor fixed, and changing the execution stack underneath them: the first run
reports \texttt{RTX 3090, 49152~MiB} with vLLM 0.27.1, the second
\texttt{RTX 4080, 32760~MiB} with vLLM 0.30.0. The weights are the same files in both
(13 files, $6{,}114{,}568{,}840$ bytes, identical per-file digests; only the directory
name differs). The outcome side reproduces exactly: all $23$ (arm, budget) cells have the
same multiset of successes, no episode changes success status ($0/672$), the monotone form's
failure to fit recurs with the same fit quality ($r^2{=}0.291377$ on both), and the
identifiability verdict is unchanged (three gains withheld against the $0.30$ CI-width
threshold). Of the $119$ leaf fields in the derived-statistics output, $117$ are equal;
the two that differ are the input path and the input digest. So Table~\ref{tab:h4} and the
behavioural mechanism it exposes do not depend, within the $8$B arm, on the first device or
the first serving
library.

Two limits on that sentence, both of which we would rather state than have a reviewer
find. First, the \emph{fitted gains} match to all printed digits, but that is structural
rather than evidential: the fit consumes per-cell success rates and nothing else, and
those are exactly what we showed to be identical. The stack change did perturb the
trajectories---$7$ of $672$ episodes used a different number of steps and $3$ a different
number of tool calls---but none of those fields reaches the fit, so agreement of
$\hat{g}_x$ across stacks adds nothing to the agreement of the rates, and we do not
describe the second run as a replication of the gains. Second, device and serving library
changed together, so this shows the finding survives a stack change, not that the device
is irrelevant; and we report no throughput comparison at all, because each run's own
device-memory readings are mutually inconsistent ($452$ versus $48{,}136$~MiB free on the
first, $4{,}102$ versus $32{,}230$~MiB on the second), so per-episode latencies from the
two runs are not on a common footing---consistent with the co-tenant caveat above.

\paragraph{The second scale: a $14$B arm that does not agree.}
The protocol registers two model arms, and both were run. The sweep above was repeated with
Qwen3-14B-AWQ, again $672$ episodes, under identical harness code (both arms record
\texttt{script\_sha256} \texttt{3326b513fa89\ldots}), the same device (\texttt{RTX 4080,
32760~MiB}), the same seed \texttt{20260911}, the same eight tasks, the same two menus and
the same five levels; the released read-out checks that identity instead of asserting it
($672/672$ records, $224/224$ \texttt{(task, menu, level)} cells and $24/24$ seeds
identical), so the only thing this pair varies is model size. Table~\ref{tab:h4-14b} gives
its curve, and it is not the curve of Table~\ref{tab:h4}. Success peaks at the \emph{mid}
level, not at the low one. The paired low-to-max contrast has \emph{more improvements than
worsenings} on the pooled records ($9$ improve, $6$ worsen, $33$ tie, $p{=}0.6072$). The
$48$ rows per level are $16$ independent (task, menu) units times three deterministic
repeats, so the paired contrasts are read on units, exactly as the $8$B arm's were:
low-to-max gives $3$ improve and $2$ worsen at $14$B ($p{=}1.00$) against $0$ improve and
$3$ worsen at $8$B ($p{=}0.25$). The behaviour offered as the mechanism reverses with
it: tests per episode \emph{rise} with the budget, $1.19\to2.00$, while rewrites rise
$1.00\to3.44$, so the verification collapse of Table~\ref{tab:h4} is not what this task set
does at $14$B. At the top budget $14$ of $16$ units pass the hidden suites against $12$ of
$16$ at $8$B ($42$/$48$ and $36$/$48$ records), and $2$ units exhaust the step budget against
$4$. The one $14$B contrast with more worsenings than improvements, mid-to-high, is $2$
worsenings and $0$ improvements on units ($p{=}0.50$; its record-level form is $6$/$0$/$42$,
$p{=}0.03125$, which we do not use)---and the $8$B arm's mid-to-high contrast points the same
way ($1$ worsening, $0$ improvements, $p{=}1.00$). The two scales therefore do not disagree
about \emph{where} the curve turns over in any statistically established sense: no paired
contrast at either scale reaches significance on $16$ units. What does differ is the location
of the peak in the descriptive curve (the low level at $8$B, the mid level at $14$B) and the
direction of the behavioural counts above.

\begin{table}[t]
\caption{The second scale: the same sweep re-run with Qwen3-14B-AWQ, $672$ episodes, under
identical code, device, seed list, tasks, menus and budget levels
(\texttt{results/h4\_derived\_stats\_14b.json}; identity check
\texttt{evidence/H4\_SCALE\_LEG\_READ.txt}). Columns as in Table~\ref{tab:h4}; $n$ is again
the pooled count of replication episodes per level. The peak is at the \emph{mid} level and
sits on the $1.0000$ ceiling, so this curve is saturated and cannot be read as an interior
optimum without that caveat. As in Table~\ref{tab:h4}, the only failure reason recorded at
any level is step-budget exhaustion ($9$ of $48$ at low, none at mid, $6$ at high, $6$ at
max); no other reason appears in the artifact.}
\label{tab:h4-14b}
\begin{center}
\small
\begin{tabular}{lrrrrr}
\toprule
Budget level & max steps & $n$ & success & writes/ep & tests/ep \\
\midrule
minimal & 2  & 48 & 0.0000 & 1.00 & 0.00 \\
low     & 4  & 48 & 0.8125 & 1.00 & 1.19 \\
mid     & 8  & 48 & \textbf{1.0000} & 1.00 & 1.25 \\
high    & 16 & 48 & 0.8750 & 1.81 & 1.69 \\
max     & 32 & 48 & 0.8750 & 3.44 & 2.00 \\
\bottomrule
\end{tabular}
\end{center}
\end{table}

\paragraph{What the $14$B arm takes back.}
The paragraph headed ``A correction we had to make to ourselves'' concluded that the
defensible statement is a
monotone decline above the low budget, \emph{not} an interior optimum, because excluding the
structurally floored level leaves the quadratic term insignificant ($p{=}0.66$, coefficient
$+0.1168$). At $14$B that same test is significant and of the opposite sign ($\chi^2{=}4.4378$
on $1$~df, $p{=}0.0352$, coefficient $-0.6375$): there the supported reading \emph{is} an
interior optimum at mid. The cited artifact disagrees with itself here, and we name it rather
than pick a side silently: \texttt{results/h4\_derived\_stats\_14b.json} sets
\texttt{non\_monotonicity\_verdict.without\_floored\_level\_quadratic\_significant} to
\texttt{true} alongside that $p{=}0.0352$, while the \texttt{reading} string in the same object
asserts that ``once excluded the quadratic is NOT significant''. We follow the numeric fields,
which are recomputable from the same file's log-likelihoods ($-66.2315$ vs $-64.0126$), not the
prose string. Three limits keep this from being a clean counter-result. (i) The
$14$B peak sits at $1.0000$, so the drop from peak to max budget is $0.1250$ of a saturated
cell and the shape comparison is confounded by ceiling: what may be said is that the $8$B
decline is not reproduced, not that the phenomenon is absent. (ii) The $p{=}0.6072$ above is
a \emph{record}-level McNemar over $48$ pairs---the same pseudo-replication that takes the
$8$B contrast from $p{=}0.0039$ to $p{=}0.25$ on units. The released \texttt{derived\_stats}
file carries unit-level recounts for the $8$B artifacts only, so we added one
(\texttt{results/h4\_unit\_recount.json}, produced by \texttt{scripts/h4\_unit\_recount.py}
from the raw \texttt{arm=iso} records, $240$ per arm, $0/16$ units varying across repeats):
low-to-max is
$3$ improve and $2$ worsen on the $16$ units ($p{=}1.00$) and mid-to-high is $2$ worsenings
and $0$ improvements ($p{=}0.50$). Neither apparent direction survives the unit recount, and
we claim neither. (iii) Two sizes of
one model family on one eight-task set are two points on a scale axis, not a scale trend.

What survives at both scales is the structural claim: the assumed saturating-monotone form is
contradicted by the measured curve at $8$B and at $14$B alike, and no fitted gain is quotable
at either scale---all three $g_x$ remain withheld against the same $0.30$ CI-width threshold
in the $14$B derivation as in the $8$B one. What does not survive is the \emph{direction} of
the low-to-max contrast and the verification-collapse mechanism proposed for it. This is not
a robustness check on Table~\ref{tab:h4}: it is a second measurement that contradicts the
claim built on the first, which is why the $8$B result is stated above as a description of a
measured curve and not as an established decline.

\section{Discussion}
\label{sec:discussion}

\paragraph{Preregistered target (H4).} H4 was preregistered to measure the three agent-axis gains
$g_{\mathrm{steps}}/g_{\mathrm{tools}}/g_{\mathrm{vision}}$ on a real model and thereby decide whether
the agent sits in a compute-limited or an accuracy-limited regime, which would predict whether the
unified router of Sec.~\ref{sec:e2} still wins. The preregistered falsifier was twofold: the gains are
quotable only if each axis has at least two interior cells \emph{and} a bootstrap CI95 width
at most $0.30$, and the joint form of Eq.~\eqref{eq:joint} requires success to be monotone
non-decreasing in every axis. Either condition failing rejects the use of the fitted gains for that
classification, not the idea of routing itself.

The results of Sec.~\ref{sec:results}, and in particular Sec.~\ref{sec:res-e3}, shape the following discussion. The source report shows that a scalar effort level gives smooth, interpolable
control over reasoning length and accuracy, and that the control transfers to
agentic trajectories (lines 2437--2443). Our allocation study suggests where that
transfer stops being free. A trajectory is not just more reasoning tokens: it is
more turns, each of which re-encodes context and may invoke a tool, so extending
the trajectory multiplies a cost that is already the dominant term. In the
token-cost optimum, reasoning saturates at the smallest budgets while trajectory
length stays at one turn, and turns are bought only after the reasoning axis is
exhausted. In latency currency the same logic pushes network search out of the
allocation entirely, because a single search call costs more wall clock
($2233.75$~ms) than an entire local tool budget.

Second, the measured dispersion argues against uniform effort tiers. The
primitives we measured span $4.8$ orders of magnitude---a $6.7\times10^{4}$ ratio between
the fastest ($0.033$~ms) and slowest ($2233.75$~ms) primitive. A deployment that raises ``effort'' by scaling every axis
together pays for the expensive axes at the same rate as the cheap ones, which is
exactly what the max-everything baseline does and why it needs $8.67\times$ the
router's tokens to reach $90\%$ of attainable success.

Third, and most consequentially, putting a real language model in the loop does
not reproduce the saturation shape (Sec.~\ref{sec:res-h4})---though the two registered
arms do not agree about what happens instead: the $8$B point estimates peak at the low
budget and fall, while the $14$B curve peaks at the mid budget and its paired low-to-max
contrast has more improvements than worsenings. This changes what the allocation study can claim.
Its frontier is conditional on the four gains being positive and saturating;
measured over one task set at two model sizes, that condition is not met, not because the
gains are small but because neither measured curve has the assumed shape. The $8$B reversal
is the striking one, and we state it as an observation rather than as an established
decline: on the $16$ independent units that contrast is $3$ worse and $0$ better,
$p{=}0.25$. We read the
mechanism as the interesting part. On the $8$B arm, effort is not a dial that scales a fixed
behaviour: a larger advertised budget changed \emph{what the agent did}, shifting
work from verification toward rewriting until the trajectory exhausted its budget
without ever producing the grading signal. That is a failure mode a cost model
over axis counts cannot express, because it depends on the policy's response to
the budget rather than on the budget itself. It also supplies a concrete instance
of the gap that \citet{wu2026ecoagentbench} report between completing a task under
a budget and spending that budget economically: here the agent is told the budget
and spends it worse. Both sentences describe the $8$B arm. The $14$B arm tested the
same mechanism on the same tasks and found tests per episode \emph{rising} with the
budget, so the mechanism is a property of that pairing, not of agent effort, and we do
not offer it as an explanation of anything the $14$B data do not show.

The source report's front-loading result and ours are not in conflict. The
report varies \emph{reasoning effort} with a trained penalty term and observes
diminishing returns; we vary \emph{agent budget} on models that were not trained
for it, and get curves with neither the assumed shape nor a common alternative one---a
reversal downward at $8$B, a peak at mid at $14$B. That distinction is itself the finding worth
stating: controllable reasoning effort does not make controllable agent effort
automatic, and the transfer the report demonstrates should not be assumed to hold
for budgets the model was never optimised against.

\paragraph{Constraints: $n$, single-axis noise, ceiling.} Three limits keep the negative result honest
rather than strong. First, the paired decline is underpowered: at $n{=}16$ independent
(task, schedule) units with paired $\alpha{=}0.05$ and power $0.8$, only standardized differences
around $0.75$ are detectable, so the observed $3$-vs-$0$ contrast with $p{=}0.25$ is directional,
not established, and not evidence of absence. The record-level $9$-vs-$0$ figure counts each of the
three deterministic repeats separately; temperature $0$ with per-request seeds leaves success unchanged
across repeats in all $80$ iso units, so the unit count is the defensible one. Second, the gain fit is
resolution-limited, not refuted in magnitude: the three CI widths are $1.00$/$1.00$/$0.58$ against the
$0.30$ bar, and each single-axis cell has only $n{=}24$ episodes, so neighbouring levels differ by
three episodes inside noise. Third, the cross-scale comparison is ceiling-confounded: the $14$B curve
peaks at $1.0000$, so its drop to $0.8750$ is $0.1250$ off a saturated cell, and two sizes of one model
family on one eight-task set are two points, not a scale trend. Wall-clock is reported only as an upper
bound under co-tenant contention and is never compared across cards.

\paragraph{Future directions.} The preregistered question stays open and now has a tighter protocol:
(i) stochastic sampling with true independent repeats so the paired test is powered at the unit level
rather than the record level; (ii) harder, unsaturated tasks where the mid budget does not sit at
$1.0000$, so shape and mechanism can be separated from ceiling; (iii) a retrieval-bearing task set
where $g_{\mathrm{search}}$ varies at all, since the current set leaves it unmeasured; (iv) a third
model scale run under the same harness and seed list, with the unit-level recount preregistered for
both arms before running. Until then the routing frontier of Sec.~\ref{sec:e2} remains conditional,
and no $g_x$ value from Sec.~\ref{sec:res-h4} may be quoted.

\paragraph{What would falsify the allocation claim.}
The claim that survives this study is comparative and conditional: \emph{if} the
agent axes have non-negligible, saturating gains, \emph{then} trading budget
across axes beats both uniform scaling and hard partition, and the optimal mix
rotates. The $30$-case grid tells us exactly which direction breaks it: weak
agent-axis gains with far-out saturation move the problem into the
accuracy-limited regime, where the reasoning axis is the only lever left. We did
measure them on real models at two scales (Sec.~\ref{sec:res-h4}), and the measurement did not
settle the regime: it returned intervals too wide to pin the gains at either scale and
curves the monotone form cannot represent. What it leaves is a concern rather than a
result: at $8$B the point estimates fall past the low setting, but on the $16$ independent
units that contrast is $3$ worse and $0$ better, $p{=}0.25$, and at $14$B the same contrast
points the other way, $3$ better and $2$ worse on the same $16$-unit design ($p{=}1.00$; its
record-level form is $9$/$6$/$33$, $p{=}0.6072$), so neither arm establishes a turnover. So the
conditional claim is untested on a real agent, and the two model-backed readings of it
disagree about the sign of the risk.
No claim in Sec.~\ref{sec:e2} depends on the assumed gains holding at any
particular value, and none of the measured $g_x$ values is quoted.

\section{Conclusion}

The discussion of Sec.~\ref{sec:discussion} closes on the same scope limit. We took the one controllable effort axis of a published systems report---a scalar
reasoning-effort level with an exponentially decaying token penalty---measured the
real cost of eighteen tool primitives on the host that runs our agent loop, and
asked what a joint router over five effort axes buys relative to uniform, fixed,
single-axis and partitioned budgets. Measured costs span $4.8$ orders of
magnitude (a $6.7\times10^{4}$ ratio); the joint router needs $3.44\times$ fewer tokens than a fixed mid
setting and $4.52\times$ fewer than maxing reasoning only at matched success,
while partitioned budgets cost $1.32\times$ more; the optimal composition rotates
across the budget rather than scaling uniformly; and an end-to-end loop with real
auto-graded tests shows monotone saturating returns with no benefit from four
times the saturated budget. We then replaced the rule-based policy with a real
language model at two scales, and the assumed shape did not survive at either: the
$8$B point estimates peak at the low budget and fall thereafter---a decline not
established on the $16$ independent units, where the contrast is $3$ worse and $0$
better, $p{=}0.25$---while the $14$B curve peaks at the mid budget and its low-to-max
contrast has more improvements than worsenings. The fitted gains were too weakly
identified to quote at either scale, and the behaviour we read as the mechanism on the
$8$B arm, a shift from verification to rewriting as the budget grew, is absent on the
$14$B arm, where test invocations per episode rise with the budget.
The honest summary is therefore two-sided. Joint routing is worth doing when the
agent axes saturate; on neither of the two real models we tested does the measured curve
have that shape, and the two disagree about how it fails, so the routing argument cannot
be extended to a real agent by assumption. We state the scope limit plainly: nothing here is
a measurement of the source model, and the four agent-axis gains remain assumptions
with a published sensitivity grid---now with two model-backed curves attached, and they
do not agree with each other.

\subsubsection*{Broader Impact Statement}
This work studies the cost side of agent deployment. Better effort routing lowers
the compute and energy cost of running an agent for the same task success, which
is a mild positive benefit. The dual risk is familiar from efficiency work:
cheaper agentic capability also lowers the barrier to running agents at scale,
including uses that are undesirable. That is a limitation of the broader effort
literature and not something a single measurement can remove; what we can do is
avoid overstating the effect. We therefore deliberately publish the full
sensitivity grid, including the parameter regimes in which our proposed router
does not help, and the loss mechanism that produces them, so that the result is
not used as a blanket claim that joint routing always wins. The tools measured
here are ordinary filesystem, process and network primitives; we performed no
security-sensitive probing, and the unreachable-host measurements target
RFC~2606 reserved names only.

\subsubsection*{Reproducibility}
All raw outputs (per-invocation latency samples, the $8{,}820$-allocation sweep,
the $240$ rule-based loop episodes, and the two $672$-episode real-model arms, $8$B and
$14$B, with their
per-episode records), the generator scripts, and the recomputation scripts are
released with the paper. Secs.~\ref{sec:e1}--\ref{sec:e3} are recomputed by
\texttt{analysis.py}, which asserts that the stored aggregates match a
recomputation from the raw samples. Sec.~\ref{sec:res-h4} is recomputed by
\texttt{h4\_derived\_stats.py} (derived statistics: the budget curve, the paired
test, both non-monotonicity fits and the per-level action profile) and
\texttt{analyze\_h4\_result.py} (\texttt{--refit} re-runs the gain fit under the
current identifiability guard). Absolute host paths are replaced by declared
placeholders in the released copies; the measured values are unchanged. Latency
measurements are host-specific and are reported with their host descriptor.

\section{Ethics and Broader Impact}

Effort routing is offered as lowering the compute and energy cost of agents, so each saving can argue for a cut made without measuring. The quotable line is that at matched success the router needs $3.44\times$ fewer tokens than a fixed mid setting. That frontier is conditional on the four agent-axis gains being positive and saturating. On the two real models we tested, over one task set, neither measured curve has that shape. The $8$B point estimates fall past a low setting rather than buying slack, though on the sixteen independent units of that arm the decline is not established, and the $14$B arm does not reproduce the fall: its curve peaks at the mid budget. Nothing in this paper is a measurement of the source model. The four agent-axis gains remain assumptions with a published sensitivity grid. The routing argument cannot be extended to a real agent by assumption, so nothing here warrants a deployed policy. The premise sits on a technical report whose own checkpoints instantiate its reported curves. Costs are host-bound: we measured eighteen tool primitives on the host that runs our loop. A search call costs more wall clock than an entire local tool budget. That $315\times$ gap is a property of this host, not a portable ratio. The warning against uniform tiers matters: scaling every axis together pays for the expensive axes at the cheap rate. That is why the max-everything baseline needs $8.67\times$ the router's tokens for $90\%$ of attainable success.

%% PROVENANCE
%% S1: Better effort routing lowers the compute and energy cost
%% S2: the joint router needs $3.44\times$ fewer tokens than a fixed mid setting
%% S3: Its frontier is conditional on the four gains being positive and saturating
%% S4: on the two real models we tested, over one task set, neither measured curve has that shape
%% S5: the $8$B point estimates fall past a low setting rather than buying slack
%% S6: nothing here is a measurement of the source model
%% S7: the four agent-axis gains remain assumptions with a published sensitivity grid
%% S8: the routing argument cannot be extended to a real agent by assumption
%% S9: only the report's own checkpoints instantiate the reported curves
%% S10: measured the real cost of eighteen tool primitives on the host
%% S11: costs more wall clock ($2233.75$~ms) than an entire local tool budget
%% S12: is $315\times$
%% S13: by scaling every axis together pays for the expensive axes
%% S14: it needs $8.67\times$ the router's tokens to reach $90\%$ of

\bibliography{main}
\bibliographystyle{tmlr}

\appendix

\section{Source quotations with line numbers}
\label{app:quotes}

This appendix lists the source quotations (App.~\ref{app:quotes}) that the paper relies on. Line numbers refer to the plain-text rendering of the source report
(SHA-256 \texttt{ac24848a\ldots abf3}), split on LF only. The extractor
\texttt{experiments/extract\_v41\_effort\_evidence.py} re-derives each quotation
and fails if a required substring is absent;
\texttt{evidence/v41flash\_effort\_control.json} holds the $40$ anchored
quotations and $19$ numeric facts.

\begin{itemize}
\item lines 2039--2044: \S5.1.4 heading and the definition of the scalar effort
  level as an explicit conditioning signal applied to both single-turn reasoning
  and multi-turn agentic tasks.
\item line 2051: ``Here, $\eff \in \{1, \ldots, 100\}$ denotes the requested
  effort level.''
\item line 2052: $M_b$ samples per prompt per effort level; lines 2061--2063:
  subgroup normalisation, ``responses from different effort levels are not
  directly compared''.
\item lines 2065--2081: the length penalty of Eq.~(9), with $\ell$ the number of
  reasoning tokens, $L_{\mathrm{norm}}$ a reference length and $C_{\max}$ the cap.
\item lines 2081--2096: exponential penalty coefficient (Eq.~10), with
  $\tau=\lambda\Delta_b$; increasing effort by $\tau$ multiplies the coefficient
  by $e^{-1}$.
\item lines 2100--2104: deployment-time interpolation between trained effort
  levels.
\item lines 2105--2109: production API tiers \texttt{max}/\texttt{high}/\texttt{low}
  at $\eff=100/75/50$; the same mapping is restated at line 2455.
\item lines 2431--2436: accuracy and output length increase with effort across
  reasoning and agentic benchmarks; $67.1\%\to76.3\%$, $66.0\%\to74.2\%$,
  $82.4\%\to90.6\%$, at roughly $2.5\times$ output tokens.
\item lines 2437--2438: ``the effort control learned on single-response reasoning
  transfers faithfully to long-horizon agentic trajectories''.
\item lines 2439--2441: front-loaded gains; $60$--$80$ recovers most of the
  maximum at less than half its token budget; the final step to $100$ lengthens
  trajectories by $1.6$--$1.8\times$.
\item lines 2456--2465: scaffold-robustness study over eight configurations from
  six scaffold families at max effort.
\item lines 2736--2763: multi-agent test-time scaling against per-rollout
  wall-clock deadlines.
\item lines 2740--2744: derived-latency accounting from token counts at fixed
  prefill/decode rates plus measured tool-execution time.
\item lines 3486--3511: appendix C, restating the length deduction (Eq.~11) and
  $k(\eff)=k_0\exp\left(-(\eff-\eff_{\min})/\tau\right)$ (Eq.~12).
\item lines 3743--3773: marginal-benefit model giving the affine effort--length
  trend (Eqs.~15--17).
\item lines 3779--3788: the explicit statement that the derivation is a local
  reward-level approximation and that realized behaviour may deviate.
\item line 24: checkpoints ``are available at'' a model hub (quoted so that the
  provenance of the source's own claims is visible).
\end{itemize}

\section{Model check: parameter calibration}
\label{app:calibration}

The reasoning kernel of Sec.~\ref{sec:kernel} is calibrated so that
(i) the accuracy endpoints reproduce lines 2434; (ii) the length ratio at
endpoints is exactly $2.5$ (line 2436); and (iii) the fraction of the
$25\to100$ accuracy gain recovered at $\reff=80$ is $0.735$, matching the
front-load statement of lines 2439--2440. The fitted curvature is $s_0=5000$ and
the fitted decay is $\kappa=0.01222$. Conditions (i)--(iii) fix both parameters
and leave no freedom; a different functional family with the same three
constraints would shift the absolute values reported in Table~\ref{tab:e2} while
leaving the ordering of the routers unchanged, which is why the ordering, not the
level, is the claim.

\end{document}
